\documentclass[12pt]{article}

% ---------------------------------------------------------
% PACKAGES
% ---------------------------------------------------------
\usepackage[a4paper,margin=0.8in]{geometry}
\usepackage{hyperref}
\usepackage{titlesec}
\usepackage{setspace}
\usepackage{enumitem}
\usepackage{microtype}
\usepackage{fancyhdr}
\usepackage{tabularx}      % For flexible-width tables
\usepackage{newtxtext,newtxmath}   % Professional Times-style font

% ---------------------------------------------------------
% DOCUMENT SETTINGS
% ---------------------------------------------------------
\hypersetup{
	colorlinks=true,
	linkcolor=black,
	urlcolor=blue,
	citecolor=black,
	pdfauthor={Leonard Okyere Afeke},
	pdftitle={MongoBleed Cyber Threat Intelligence Report}
}

\setstretch{1.2}

\titleformat{\section}{\large\bfseries}{\thesection}{1em}{}
\titleformat{\subsection}{\normalsize\bfseries}{\thesubsection}{1em}{}

\pagestyle{fancy}
\fancyhf{}
\rhead{MongoBleed CTI Report}
\lhead{Leonard Okyere Afeke}
\cfoot{\thepage}

% ---------------------------------------------------------
% TITLE PAGE
% ---------------------------------------------------------
\begin{document}
	
	\begin{titlepage}
		\centering
		\vspace*{3cm}
		{\Huge \textbf{MongoBleed}}\par
		\vspace{0.5cm}
		{\Large Cyber Threat Intelligence Case Study}\par
		\vspace{1.5cm}
		
		{\large \textbf{Author:} Leonard Okyere Afeke}\par
		\vspace{0.3cm}
		{\large \textbf{Date:} \today}\par
		\vspace{2cm}
		
		\vfill
		{\small This report is based entirely on publicly available information (OSINT). 
			No scanning, exploitation, or interaction with live systems occurred.}
	\end{titlepage}
	
	% ---------------------------------------------------------
	% EXECUTIVE SUMMARY
	% ---------------------------------------------------------
	\section*{Executive Summary}
	
	MongoBleed is a long‑running pattern of opportunistic attacks targeting \textbf{publicly exposed MongoDB instances} left accessible without authentication.  
	Attackers rapidly scan the internet for open databases, delete stored data, and leave ransom notes claiming the data has been “backed up.” Public reporting consistently shows these claims are false — the attacks result in \textbf{irreversible data loss}.
	
	This report synthesizes publicly available information to provide a structured Cyber Threat Intelligence (CTI) assessment of MongoBleed, including attacker behavior, infrastructure patterns, ransom note characteristics, and defensive recommendations.
	
	MongoBleed demonstrates the ongoing risks of cloud misconfiguration and highlights the need for continuous exposure monitoring, secure defaults, and strong access controls.
	
	% ---------------------------------------------------------
	% KEY JUDGMENTS
	% ---------------------------------------------------------
	\section{Key Judgments}
	
	\begin{itemize}
		\item MongoBleed is \textbf{not a single threat actor}, but a recurring pattern of opportunistic exploitation performed by multiple low‑sophistication groups.
		\item Attackers rely heavily on \textbf{automated scanning} to identify exposed MongoDB instances within minutes of exposure.
		\item Infrastructure used in these attacks is typically \textbf{short‑lived VPS servers}, often clustered around specific ASNs.
		\item Ransom notes follow \textbf{highly consistent templates}, with repeated claims of data backup and Bitcoin payment demands.
		\item Public reporting shows \textbf{no evidence} that attackers store or back up victim data; restoration is effectively impossible.
		\item MongoBleed attacks are \textbf{destructive extortion}, not true ransomware.
		\item The threat persists due to \textbf{misconfigurations}, not software vulnerabilities.
	\end{itemize}
	
	% ---------------------------------------------------------
	% THREAT OVERVIEW
	% ---------------------------------------------------------
	\section{Threat Overview}
	
	MongoBleed occurs when MongoDB instances are deployed:
	
	\begin{itemize}
		\item without authentication  
		\item exposed directly to the internet  
		\item using legacy or insecure defaults  
		\item without firewall or security group restrictions  
	\end{itemize}
	
	Attackers exploit these conditions by:
	
	\begin{itemize}
		\item scanning TCP/27017  
		\item connecting anonymously  
		\item enumerating databases  
		\item deleting data  
		\item inserting ransom notes  
	\end{itemize}
	
	This attack chain is simple, automated, and highly scalable.
	
	% ---------------------------------------------------------
	% OBSERVED ATTACKER BEHAVIOR
	% ---------------------------------------------------------
	\section{Observed Attacker Behavior}
	
	\subsection*{Discovery}
	Public reporting (e.g., Rapid7, Shodan) shows attackers use:
	
	\begin{itemize}
		\item automated internet‑wide scanning  
		\item search engines indexing exposed services  
		\item scripts that test for unauthenticated access  
	\end{itemize}
	
	\subsection*{Exploitation}
	Once an exposed instance is found:
	
	\begin{itemize}
		\item attackers connect without credentials  
		\item enumerate databases  
		\item delete collections  
		\item create ransom note collections  
	\end{itemize}
	
	\subsection*{Post‑Exploitation}
	Attackers typically:
	
	\begin{itemize}
		\item do not maintain persistence  
		\item do not monitor payments  
		\item do not store victim data  
	\end{itemize}
	
	MongoBleed campaigns are \textbf{short‑cycle and disposable}.
	
	% ---------------------------------------------------------
	% RANSOM NOTE CHARACTERISTICS
	% ---------------------------------------------------------
	\section{Ransom Note Characteristics}
	
	Public OSINT sources show that ransom notes follow a consistent structure.
	
	\subsection*{Common Phrases (short excerpts)}
	\begin{quote}
		“Your database has been backed up…”\\
		“Send X BTC to the following address…”
	\end{quote}
	
	\subsection*{Paraphrased Summary}
	Ransom notes typically:
	
	\begin{itemize}
		\item claim the attacker exported the victim’s data  
		\item demand Bitcoin for restoration  
		\item provide a wallet and contact method  
		\item threaten permanent data loss  
	\end{itemize}
	
	\subsection*{Reconstructed Template}
	
	\begin{verbatim}
		READ_ME
		--------------------
		Your database has been backed up.
		To recover your data, send X BTC to [WALLET_ADDRESS].
		After payment, contact us at [EMAIL/TELEGRAM].
		If you do not pay, your data will be permanently deleted.
	\end{verbatim}
	
	\subsection*{Key Insight}
	Public reporting (e.g., BleepingComputer, Rapid7) shows \textbf{no evidence} that attackers store or back up data.  
	Restoration is impossible.
	
	% ---------------------------------------------------------
	% INFRASTRUCTURE ANALYSIS
	% ---------------------------------------------------------
	\section{Infrastructure Analysis}
	
	\subsection*{VPS Usage}
	Attackers frequently use VPS providers such as:
	
	\begin{itemize}
		\item DigitalOcean  
		\item OVH  
		\item Hetzner  
		\item Linode  
	\end{itemize}
	
	These providers enable:
	
	\begin{itemize}
		\item low‑cost automation  
		\item rapid provisioning  
		\item disposable infrastructure  
	\end{itemize}
	
	\subsection*{ASN Clustering}
	Scanning IPs often cluster around specific ASNs, indicating:
	
	\begin{itemize}
		\item repeated use of the same hosting provider  
		\item automated deployment across multiple IPs  
	\end{itemize}
	
	\subsection*{Geographic Distribution}
	Infrastructure is globally distributed, reflecting hosting provider availability rather than attacker origin.
	
	% ---------------------------------------------------------
	% MITRE ATT&CK MAPPING
	% ---------------------------------------------------------
	\section{MITRE ATT\&CK Mapping}
	
	\begin{table}[h!]
		\centering
		\begin{tabularx}{\textwidth}{|l|l|X|}
			\hline
			\textbf{Phase} & \textbf{Technique} & \textbf{Description} \\
			\hline
			Reconnaissance & Active Scanning (T1595) & Automated scanning of TCP/27017 \\
			\hline
			Resource Development & Acquire Infrastructure (T1583.003) & Use of VPS providers \\
			\hline
			Initial Access & Exploit Public‑Facing Application (T1190) & Accessing unauthenticated MongoDB \\
			\hline
			Discovery & Query Databases (T1505.003) & Enumerating databases and collections \\
			\hline
			Impact & Data Destruction (T1485) & Dropping databases \\
			\hline
			Impact & Defacement (T1491) & Creating ransom note collections \\
			\hline
		\end{tabularx}
	\end{table}
	
	% ---------------------------------------------------------
	% DEFENSIVE RECOMMENDATIONS
	% ---------------------------------------------------------
	\section{Defensive Recommendations}
	
	\subsection*{1. Preventive Controls}
	\begin{itemize}
		\item Enable MongoDB authentication  
		\item Bind MongoDB to localhost or private networks  
		\item Restrict access using firewalls or cloud security groups  
		\item Disable legacy configuration defaults  
		\item Use managed database services with secure defaults  
	\end{itemize}
	
	\subsection*{2. Detection Opportunities}
	\begin{itemize}
		\item Anonymous access attempts  
		\item Unexpected connections to TCP/27017  
		\item Rapid database deletion events  
		\item Creation of suspicious collections (e.g., “READ\_ME”)  
		\item Connections from VPS‑associated ASNs  
	\end{itemize}
	
	\subsection*{3. Hardening Measures}
	\begin{itemize}
		\item Enforce role‑based access control  
		\item Enable auditing and logging  
		\item Require TLS for all connections  
		\item Continuously scan cloud assets for misconfigurations  
	\end{itemize}
	
	% ---------------------------------------------------------
	% CONCLUSION
	% ---------------------------------------------------------
	\section{Conclusion}
	
	MongoBleed is a persistent, automated threat driven by cloud misconfigurations rather than sophisticated adversaries.  
	Its impact is severe due to irreversible data loss, rapid exploitation, and widespread exposure of unsecured databases.
	
	Organizations can significantly reduce risk by:
	
	\begin{itemize}
		\item enforcing authentication  
		\item restricting network exposure  
		\item adopting secure deployment practices  
		\item monitoring for misconfigurations  
	\end{itemize}
	
	MongoBleed underscores the importance of \textbf{cloud hygiene}, \textbf{secure defaults}, and \textbf{continuous exposure management}.
	
	% ---------------------------------------------------------
	% SOURCES
	% ---------------------------------------------------------
	\section*{Sources Consulted (Public OSINT)}
	
	\begin{itemize}
		\item BleepingComputer — reporting on MongoDB ransom attacks  
		\item Rapid7 — research on exposed databases and opportunistic exploitation  
		\item Shodan Blog — analysis of exposed services and scanning behavior  
		\item Censys — internet‑wide exposure insights  
		\item MongoDB Documentation — security best practices and authentication guidance  
		\item Public GitHub repositories — examples of MongoDB ransom scripts  
		\item Comparitech — exposed database research  
	\end{itemize}
	
\end{document}
